% 06-experiences-resultats (fr): every experiment's result, ten seeds first, then twenty seeds and equal budget.
\chapter{Les expériences : résultats}\label{chap:06-experiences-resultats}

Ce chapitre donne le résultat de chaque expérience décrite au \cref{chap:05-experiences-methode}. Chaque
section suit le même ordre : la question, le plan en une phrase, le résultat avec ses intervalles, puis ce
qu'il signifie et ce qu'il ne dit pas.

\paragraph{Conventions de lecture}
Les résultats principaux sont ceux à dix graines aléatoires (0 à 9), annoncés par le protocole. Les
résultats à vingt graines aléatoires (0 à 19) et à budget égal viennent ensuite. Chaque divergence est dite. Un score
est le retour moyen des dix épisodes d'évaluation finaux d'un essai (\cref{sec:evaluation}). Un score de
programme est une moyenne interquartile (IQM) de ces scores sur les graines aléatoires. Un intervalle écrit
entre crochets, comme $[\num{2268}, \num{3369}]$, est un intervalle bootstrap à 95~\%. Un écart est
\emph{établi} quand il satisfait la règle du vainqueur établi (\cref{sec:regle}) : intervalles des IQM
disjoints, ou intervalle de la probabilité d'amélioration qui exclut 0,5. Les tableaux sont produits par le
script \file{report/tools/tables.py} à partir des fichiers JSON de \file{sim/analysis/}. Les nombres du texte
viennent des mêmes fichiers, avec l'arrondi de \file{docs/conclusions.md}.

\section{E1 : même code, autre graine aléatoire}\label{sec:e1}

\paragraph{Question et plan}
Deux groupes d'essais du même programme, qui ne diffèrent que par leurs graines aléatoires, peuvent-ils
passer pour deux algorithmes différents ? E1 entraîne trente fois PPO PyTorch sur \task{HalfCheetah-v5}, un
million de pas chacun, graines aléatoires 0 à 29, puis compare par un test de Welch des groupes tirés parmi
ces trente essais (\cref{sec:plan-e1}).

\paragraph{Une dispersion large, qui n'a pas la forme d'une cloche}
La moyenne des trente scores vaut \num{1688}, l'écart type 568, soit 34~\% de la moyenne. La médiane vaut
\num{1580}. Vingt-sept essais finissent entre \num{1138} et \num{1999}. Trois décollent au-delà de \num{3000},
jusqu'à \num{3430} (\cref{fig:e1-graines}). Ces trois essais tirent la moyenne au-dessus de la médiane. Le test
de Shapiro-Wilk rejette nettement la loi normale ($p = 4\times10^{-7}$). Un test t suppose des moyennes de
groupe à peu près normales : avec cinq essais tirés dans une telle distribution, l'hypothèse est fragile.

\fig{e1-graines}{Trente entraînements du même programme sur la même tâche : PPO écrit en PyTorch
(Stable-Baselines3) sur \task{HalfCheetah-v5}, un million de pas chacun. Chaque point est un essai, placé à son
score d'évaluation final, la moyenne des retours de dix épisodes joués sans hasard sur des graines aléatoires
jamais vues à l'entraînement. Seule la graine aléatoire d'entraînement change d'un point à l'autre. Moyenne
\num{1688}, médiane \num{1580} (ligne pointillée). Vingt-sept essais finissent entre \num{1138} et \num{1999},
trois au-delà de \num{3000}.}

\paragraph{Le hasard trompe rarement le test, le tri le trompe souvent}
Le \cref{tab:e1} et la \cref{fig:e1-faux-positifs} donnent la part des \num{10000} tirages où le test de Welch
déclare différents, au seuil de 5~\%, deux groupes issus du même programme. Chacune de ces déclarations est un
faux positif.
\begin{itemize}
  \item Trois essais contre trois, tirés au hasard : 1,6~\% des tirages.
  \item Cinq contre cinq : 1,2~\%.
  \item Dix contre dix : 1,7~\%.
  \item Les cinq meilleurs contre les cinq pires d'un tirage de dix essais : 31,0~\%.
\end{itemize}
Tirés au hasard, les groupes trompent le test moins souvent que le seuil nominal de 5~\%. Le test est
conservateur sur cette distribution asymétrique. L'explication la plus probable : un essai hors norme gonfle
la variance du groupe qui le reçoit en même temps que sa moyenne. Triés, les groupes le trompent près d'une fois sur trois. La valeur $p$
médiane de la coupe triée vaut 0,09.

\gentable{e1}

\fig{e1-faux-positifs}{Part des \num{10000} tirages où le test de Welch déclare différents, au seuil de 5~\%,
deux groupes d'essais du même programme (PPO PyTorch sur \task{HalfCheetah}, trente essais qui ne diffèrent
que par leur graine aléatoire). Chaque grille compte 100 cases, une case valant un centième des tirages. Les
cases encadrées en pointillé figurent le seuil de 5~\%. Tirés au hasard, trois contre trois, cinq contre cinq ou
dix contre dix, les groupes trompent le test moins souvent que ce seuil. Les cinq meilleurs opposés aux cinq
pires d'un tirage de dix le trompent près d'une fois sur trois.}

\paragraph{Le hasard trompe l'œil}
Les groupes tirés au hasard trompent rarement le test, mais leurs moyennes diffèrent beaucoup. L'écart relatif
des moyennes, rapporté à la plus petite des deux, vaut en médiane 17~\% pour trois contre trois, 17~\% pour
cinq contre cinq et 11~\% pour dix contre dix. Une fois sur vingt, il dépasse 62~\%, 47~\% et 33~\%. Deux
groupes de cinq essais tirés au hasard dessinent souvent deux courbes bien séparées. Dans la coupe triée,
l'écart médian atteint 42~\%.

\paragraph{Garder les trois meilleurs essais sur dix}
Sur \num{10000} tirages de dix essais, la moyenne des trois meilleurs vaut \num{2249} en moyenne, contre
\num{1692} pour les dix : un gonflement de 33~\%. L'expérience E4 étend cette mesure aux quatre programmes
(\cref{sec:e4}).

\paragraph{La statistique de 2018, hors de portée de PPO}
\textcite{henderson2018matters} testent des courbes d'apprentissage entières. La seconde statistique de E1,
ajoutée après les résultats précédents (écart~3), refait leur t moyen sur nos trente courbes : un test de Welch
à chacun des vingt points d'évaluation, puis la moyenne des $t$ (\cref{sec:t-moyen}). Une fois sur vingt, la
valeur absolue du t moyen dépasse 1,4 quand les groupes de cinq sont tirés au hasard, 2,5 quand ils sont triés
par score final. Aucun des \num{20000} tirages n'atteint 9. Le test au seul dernier point retrouve le résultat
des scores finaux : environ 31~\% de faux positifs pour les groupes triés.

\paragraph{Ce que E1 établit, ses limites}
Sur ces trente essais, le hasard des graines aléatoires trompe l'œil bien plus souvent que le test. Un test de
Welch sur des groupes tirés au hasard reste prudent, même avec trois essais par groupe. Le tri des essais
change ce hasard en faux positifs. La sélection des meilleurs gonfle la moyenne d'un tiers. E1 se limite à
un programme et une tâche. La forme de la distribution, avec quelques essais qui décollent, dépend des deux.
Sur une distribution plus symétrique, le test au hasard serait plus proche de son seuil nominal.

\section{B1 : l'énigme du t de 9}\label{sec:b1}

\paragraph{Question}
La figure~5 de \textcite{henderson2018matters} oppose deux groupes de cinq essais de TRPO sur
\task{HalfCheetah}, avec le même code et les mêmes réglages. Elle rapporte $t=-9{,}09$ et $p=0{,}0016$. Pineau
précise en 2018 que les groupes ne sont pas triés. Or E1 montre qu'avec PPO, des groupes tirés au hasard
n'approchent jamais un t moyen de 9. D'où vient ce chiffre ? Trois pistes restaient ouvertes : la dispersion
propre à TRPO, une autre façon de lire leur légende, ou un autre calcul (\cref{sec:plan-b1}).

\paragraph{Première piste : TRPO disperse davantage, sans approcher 9}
Les trente essais de TRPO, entraînés avec les valeurs par défaut de sb3-contrib, dispersent bien plus que ceux
de PPO. Leur moyenne finale vaut \num{1951}, leur écart type 930, soit 48~\% de la moyenne, contre 34~\% pour
PPO. Pourtant, avec des groupes de cinq tirés au hasard, le t moyen sur la courbe ne dépasse 1,4 qu'une fois
sur vingt, comme pour PPO. Avec des groupes triés, il dépasse 1,9 une fois sur vingt. La règle fixée d'avance
(écart~4) écarte la dispersion de l'algorithme comme explication.

\paragraph{Deuxième piste : les points empilés, un calcul fautif qui reste court}
La troisième lecture empile les cent points des cinq courbes de chaque groupe, comme s'ils étaient des
observations indépendantes. Une fois sur vingt, son $|t|$ dépasse 5,0 avec des groupes tirés au hasard, 6,3
avec des groupes triés. Aucun tirage n'atteint 9. Cette lecture a un défaut plus grave que son échec : elle
déclare une différence dans 53~\% des tirages au hasard avec TRPO, 23~\% avec PPO, alors que les deux groupes
font tourner le même code. Les vingt points d'une courbe ne sont pas indépendants : un essai qui part bien reste
souvent devant.

\paragraph{Le p publié trahit le calcul}
Un test de Welch bilatéral qui donne $t=9{,}09$ et $p=0{,}0016$ a environ 3,4~degrés de liberté. Ce couple
exclut les points empilés, qui ont environ 198~degrés de liberté et donneraient un $p$ de l'ordre de
$10^{-16}$. Il exclut aussi le t moyen converti en $p$ : deux groupes de cinq donnent au moins 4~degrés de
liberté (\cref{sec:welch}). Seule une moyenne de valeurs $p$ reste compatible avec le couple publié.

\paragraph{Le script de 2017}
Peter Henderson a publié en juin 2017 un script au nom parlant, \file{from\_same\_dist.py}, \enquote{de la même
distribution} \parencite{henderson2017samedist}. La \cref{sec:script-henderson} en donne le calcul exact. En
résumé :
\begin{enumerate}
  \item il fait la moyenne des courbes des cinq graines aléatoires de chaque groupe, point par point ;
  \item à chaque point, il compare les deux courbes moyennes par un test de Welch dont la variance est celle
    \emph{des épisodes} d'une même itération, divisée par le nombre d'épisodes ;
  \item il rapporte la moyenne des $t$ et la moyenne des $p$.
\end{enumerate}
La variabilité entre graines aléatoires, celle que la figure devait mesurer, n'entre jamais dans le test. Le
script compte les épisodes au lieu des entraînements indépendants.

\paragraph{Le script retrouve le t de 9 sur nos essais de TRPO}
Appliqué à nos trente courbes de TRPO avec $n=10$ épisodes, selon un plan écrit avant le calcul (écart~6), le
script donne les résultats du \cref{tab:b1} :
\begin{itemize}
  \item Avec des groupes tirés au hasard, la valeur absolue du t moyen atteint 9,09 dans 38~\% des tirages.
    Sa médiane vaut 7,1. Le test déclare une différence dans 43~\% des tirages.
  \item Avec des groupes triés, 70~\% des tirages atteignent 9,09.
  \item Avec $n=5$, l'analyse de sensibilité, 23~\% des tirages au hasard atteignent 9,09. Le test y déclare
    une différence dans 25~\% des tirages.
\end{itemize}
Appliqué à nos trente courbes de PPO, le même calcul n'atteint jamais 9,09. Une fois sur vingt, son $|t|$
moyen dépasse 3,6 au hasard. La \cref{fig:t-de-neuf} montre ce contraste, la \cref{fig:b1-distribution}
montre les distributions complètes.

\gentable{b1}

\paragraph{Pourquoi TRPO seul y arrive}
Le dénominateur du script est l'écart type entre épisodes. Le bon dénominateur est l'écart type entre graines
aléatoires. Notons $\rho$ le rapport du premier au second, mesuré au dernier point d'évaluation. Plus $\rho$
est petit, plus le $t$ du script dépasse celui d'un test sur les graines aléatoires, d'un facteur proche de $1/\rho$. Chez
notre TRPO, l'écart type entre les épisodes d'un même agent ne vaut que 7~\% de l'écart type entre agents. Chez
notre PPO, il en vaut 30~\%. Un agent TRPO joue presque le même score à chaque épisode, alors que deux agents
TRPO diffèrent beaucoup. C'est le cas où compter les épisodes au lieu des graines aléatoires fabrique des $t$
énormes.

\fig{t-de-neuf}{Deux groupes de cinq essais tirés au hasard parmi trente, \num{10000} fois, pour chaque
algorithme et chaque façon de calculer le $t$ moyen le long de la courbe d'apprentissage. Chaque grille compte
100 cases, une case valant un centième des tirages. Elle colore la part des tirages où le $t$ moyen, signe
oublié, atteint 9,09, la valeur publiée par Henderson et ses coauteurs en 2018 pour TRPO sur
\task{HalfCheetah}. Seul le script qu'ils ont publié en 2017, qui teste la dispersion entre épisodes au lieu de
la dispersion entre graines aléatoires, y parvient, appliqué à TRPO : 38~\% des tirages. PPO avec ce script,
PPO et TRPO avec notre test sur les graines aléatoires, TRPO avec les points empilés : aucun tirage.}

\fig{b1-distribution}{Distribution de la valeur absolue du $t$ moyen sur \num{10000} tirages de deux groupes de
cinq essais pris au hasard parmi trente, sur \task{HalfCheetah}, pour TRPO (en haut) et pour PPO PyTorch (en
bas), en échelle logarithmique. En gris, le test de Welch fait entre graines aléatoires à chaque point de la
courbe, puis moyenné. En rouge, le calcul du script de Henderson (2017), qui fait la moyenne des courbes de
chaque groupe et teste la dispersion entre les dix épisodes d'évaluation. La ligne verticale marque 9,09, la
valeur publiée en 2018. Seul le script appliqué à TRPO la dépasse, dans 38~\% des tirages.}

\paragraph{Ce qui reste incertain}
Nous ne pouvons pas prouver que ce script a produit la figure~5. Son nom, sa sortie (\enquote{t=\dots, p=\dots}) et sa
compatibilité avec le couple $(t,p)$ publié le rendent probable. Nos conditions diffèrent des leurs. Nous
l'avions dit avant le calcul. Nos épisodes d'évaluation sont joués par une politique déterministe, alors que
leurs fichiers de progression enregistrent les épisodes d'entraînement d'une politique stochastique, qui
varient sans doute davantage. Nous avons vingt points d'évaluation au lieu de leurs itérations. Notre TRPO est
celui de sb3-contrib. Leur code était différent. L'ordre de grandeur se retrouve malgré ces différences.

\paragraph{Ce que B1 établit}
Le t de 9 de 2018 n'exige aucun tri des essais pour apparaître : un test qui compte les épisodes au lieu des
graines aléatoires suffit à le produire sur nos propres essais de TRPO. La faille F11, même code, autre graine
aléatoire, est plus profonde que la figure ne le montrait : le hasard des graines aléatoires a trompé le test même qui devait le mesurer. La
\cref{fig:e1-henderson} rassemble E1 et B1 face au chiffre publié.

\fig{e1-henderson}{Nos trente essais de PPO et nos trente essais de TRPO sur \task{HalfCheetah}, comparés au
chiffre publié en 2018. Pour chaque façon de répartir dix essais en deux groupes de cinq, au hasard ou triés
par score final, la barre donne la valeur absolue du $t$ moyen le long de la courbe d'apprentissage dépassée une
fois sur vingt, sur \num{10000} tirages. La barre des points empilés traite les cent points des courbes de
chaque groupe comme des observations indépendantes. La barre du script de Henderson applique son calcul de
2017, qui teste la dispersion entre épisodes, à nos essais de TRPO. La barre de Henderson et ses coauteurs est
la valeur qu'ils publient sur leurs propres essais.}

\section{E2 : le programme en tête change de tâche en tâche}\label{sec:e2}

\paragraph{Question et plan}
\textcite{henderson2018matters} montraient que l'algorithme en tête sur une tâche cessait de l'être sur une
autre. E2 refait la comparaison avec quatre programmes de 2026 (PPO PyTorch, PPO JAX, SAC, CrossQ) sur les
quatre tâches, dix graines aléatoires chacun, au budget de chaque programme : un million de pas pour PPO,
\num{200000} pour SAC et CrossQ (\cref{sec:plan-e2}).

\paragraph{Dix graines aléatoires : quatre tâches, trois programmes en tête}
Le \cref{tab:e2} donne l'IQM de chaque programme sur chaque tâche avec son intervalle. La
\cref{fig:classement-taches} rapporte chaque IQM à la meilleure de sa tâche.
\begin{itemize}
  \item \task{HalfCheetah} : SAC mène avec \num{5191} $[\num{4431}, \num{5601}]$, devant CrossQ, \num{4801}
    $[\num{2427}, \num{7815}]$. Les deux PPO suivent loin derrière, \num{1407} pour PPO PyTorch et \num{1254}
    pour PPO JAX.
  \item \task{Hopper} : PPO JAX mène avec \num{3064}, devant PPO PyTorch (\num{2877}), CrossQ (\num{2302}) et
    SAC (\num{1308}).
  \item \task{Walker2d} : CrossQ mène avec \num{5414} $[\num{4994}, \num{5627}]$, loin devant PPO PyTorch
    (\num{2158}), PPO JAX (\num{1272}) et SAC (\num{1236}).
  \item \task{Swimmer} : PPO JAX mène avec 112, devant CrossQ (110) et PPO PyTorch (107). SAC reste à 47.
\end{itemize}

\fig{classement-taches}{Score de chaque programme rapporté au meilleur de sa tâche, sur dix graines
aléatoires : la barre à 100 désigne le programme en tête. Le score d'un programme est la moyenne
interquartile (la moyenne des six scores du milieu sur dix) de ses scores d'évaluation finaux. PPO s'entraîne
un million de pas, SAC et CrossQ \num{200000}. Le programme en tête change : SAC sur \task{HalfCheetah}, PPO
écrit en JAX sur \task{Hopper} et \task{Swimmer}, CrossQ sur \task{Walker2d}.}

\paragraph{Un changement établi, une tête rarement nette}
Le programme en tête n'est pas le même partout. Le changement est établi selon la règle du protocole : SAC, en
tête sur \task{HalfCheetah}, est battu de façon tranchée sur \task{Hopper} par les deux PPO, sur
\task{Walker2d} par CrossQ et PPO PyTorch, sur \task{Swimmer} par les trois autres. \file{e2.py} compte
14~changements établis. En revanche, le programme en tête ne bat de façon tranchée tous les autres que sur
\task{Walker2d}. Sur les trois autres tâches, dix graines aléatoires ne le séparent pas de son suivant : SAC de
CrossQ sur \task{HalfCheetah}, PPO JAX de PPO PyTorch sur \task{Hopper}, PPO JAX de CrossQ sur
\task{Swimmer}. L'intervalle de CrossQ sur \task{HalfCheetah}, de \num{2427} à \num{7815}, montre combien ses
essais varient sur cette tâche.

\paragraph{Avec la correction de Bonferroni}
E2 calcule 24~intervalles de probabilité d'amélioration sans correction. Portés au niveau simultané de 95~\%,
soit 99,79~\% chacun (écart~18), ils laissent 12~changements établis au lieu de 14 (\cref{tab:e2-bonferroni}).
Deux paires perdent leur statut : PPO JAX contre CrossQ sur \task{Hopper}, PPO PyTorch contre SAC sur
\task{Walker2d}. Le verdict de E2 résiste à la correction.

\gentable{e2}

\gentable{e2-bonferroni}

\paragraph{Les courbes d'apprentissage}
La \cref{fig:courbes-apprentissage} montre les courbes d'évaluation des seize couples. Elle rappelle ce qu'un
classement final cache : les courbes de SAC et de CrossQ s'arrêtent à \num{200000}~pas, celles de PPO
continuent jusqu'à un million. La lecture à budget égal, plus bas, compare les quatre programmes au même
point de ces courbes.

\fig{courbes-apprentissage}{Courbes d'évaluation des quatre programmes (PPO PyTorch, PPO JAX, SAC, CrossQ) sur
les quatre tâches (une tâche par colonne). En abscisse, les pas d'entraînement, un point tous les
\num{50000}~pas : jusqu'à un million pour PPO, \num{200000} pour SAC et CrossQ. En ordonnée, le retour moyen
de dix épisodes d'évaluation joués sans hasard. Le trait plein est la moyenne interquartile des dix graines
aléatoires, la bande son intervalle bootstrap à 95~\%, les traits pâles les essais individuels.}

\paragraph{Vingt graines aléatoires}
Avec les graines aléatoires 0 à 19, le programme en tête change toujours avec la tâche : CrossQ mène sur
\task{HalfCheetah} et \task{Walker2d}, PPO JAX sur \task{Hopper} et \task{Swimmer}. Il reste 10~changements
établis, 6 avec la correction de Bonferroni. Trois divergences avec les dix graines aléatoires sont à dire.
\begin{itemize}
  \item Sur \task{HalfCheetah}, SAC cède la tête à CrossQ : son IQM passe de \num{5191} à \num{4214}, celle de
    CrossQ de \num{4801} à \num{4694}. L'écart entre les deux n'est établi ni à dix ni à vingt graines aléatoires.
  \item Sur \task{Hopper}, CrossQ monte de la troisième à la deuxième place (\num{2302} puis \num{2792}).
    L'avance de PPO JAX sur lui n'est plus établie.
  \item Le programme en tête bat désormais de façon tranchée tous les autres sur deux tâches,
    \task{Walker2d} et \task{Swimmer}.
\end{itemize}

\paragraph{À budget égal}
PPO reçoit cinq fois plus de pas que SAC et CrossQ. Lu à \num{200000}~pas sur sa courbe (écart~12, point~iii),
il ne mène plus sur aucune tâche.
\begin{itemize}
  \item À dix graines aléatoires, SAC mène sur \task{HalfCheetah}, CrossQ sur les trois autres tâches. Il reste
    3~changements établis : SAC, en tête sur \task{HalfCheetah}, est battu de façon tranchée par CrossQ sur
    \task{Walker2d} et sur \task{Swimmer}, puis par PPO JAX sur \task{Swimmer}. Sur \task{Hopper}, aucune
    paire n'est tranchée.
  \item À vingt graines aléatoires, CrossQ mène sur les quatre tâches et aucun changement n'est établi. Sur
    \task{HalfCheetah}, son avance sur SAC (\num{4694} contre \num{4214}) n'est pas établie non plus.
\end{itemize}

\paragraph{Ce que E2 établit, ses limites}
Le verdict préenregistré tient : à dix graines aléatoires et aux budgets du protocole, le programme en tête
change de tâche en tâche, de façon établie, même après correction de Bonferroni. Deux nuances pèsent. Le
programme en tête bat rarement tous les autres de façon nette. Avec vingt graines aléatoires, le changement de tête d'une
tâche à l'autre vient en bonne part du budget cinq fois plus grand de PPO : à budget égal, CrossQ mène
partout. Un article qui compare deux méthodes à budgets différents compare aussi ces budgets. Enfin, le
classement vaut pour les valeurs par défaut des bibliothèques. Un réglage par tâche le changerait peut-être,
dans un sens que nous ignorons : l'expérience E6 montre, sur un cas, qu'un réglage publié fait moins bien que
les valeurs par défaut.

\section{E3 : deux programmes du même algorithme}\label{sec:e3}

\paragraph{Question et plan}
Deux programmes qui réalisent le même algorithme donnent-ils le même résultat ? E3 oppose PPO PyTorch à PPO
JAX sur les quatre tâches, avec les mêmes graines aléatoires, les mêmes valeurs par défaut et un million de pas
(\cref{sec:plan-e3}). Les deux programmes diffèrent par la bibliothèque de calcul et par l'initialisation des
poids (\cref{sec:programmes}).

\paragraph{Dix graines aléatoires : aucun vainqueur, des écarts larges}
Le \cref{tab:e3} donne les deux IQM, leur rapport et la probabilité qu'un essai de PPO JAX batte un essai de
PPO PyTorch. Sur aucune tâche, l'une des deux versions ne l'emporte selon la règle du vainqueur établi. Les
écarts apparents sont pourtant larges. Sur \task{Walker2d}, l'IQM vaut \num{2158} avec PyTorch et \num{1272}
avec JAX, 41~\% de moins. La probabilité que JAX fasse mieux vaut 0,27, de 0,06 à 0,52 : l'intervalle touche
encore 0,5. Les deux intervalles des IQM se chevauchent. Sur les trois autres tâches, les probabilités
d'amélioration restent proches de 0,5 : 0,36 sur \task{HalfCheetah}, 0,54 sur \task{Hopper}, 0,66 sur
\task{Swimmer}. Les intervalles des IQM de \file{e3.json} diffèrent de quelques points de ceux de
\file{e2.json} pour les mêmes essais, parce que chaque script consomme le générateur aléatoire dans son propre
ordre (\cref{sec:reproductibilite-bootstrap}).

\gentable{e3}

\paragraph{Vingt graines aléatoires : les deux versions se séparent}
C'est le résultat qui bouge le plus avec les graines aléatoires ajoutées. Avec vingt graines aléatoires, deux tâches ont un
vainqueur établi :
\begin{itemize}
  \item Sur \task{Walker2d}, PyTorch l'emporte, \num{2204} contre \num{1435}. La probabilité que JAX fasse
    mieux vaut 0,30, de 0,15 à 0,48.
  \item Sur \task{Swimmer}, JAX l'emporte, 115 contre 107. La probabilité vaut 0,79, de 0,63 à 0,91, avec des
    intervalles d'IQM disjoints.
\end{itemize}
Les écarts ressemblent à ceux de dix graines aléatoires. Les dix ajoutées resserrent les intervalles.

\paragraph{Ce que E3 établit, ses limites}
Avec dix graines aléatoires, nous ne pouvons pas dire que les deux versions de PPO diffèrent. Avec vingt, elles
diffèrent sur deux tâches, dans des sens opposés. Une méthode comparée à \enquote{PPO} sur \task{Walker2d}
gagnerait ou perdrait selon la version choisie. La faille F14, l'implémentation qui change le résultat, tient
en 2026 entre deux bibliothèques du même auteur, aux valeurs par défaut identiques. E3 ne dit pas laquelle des
différences de code en est la cause : l'initialisation des poids, les générateurs aléatoires ou les arrondis de
la bibliothèque de calcul.

\section{E4 : les trois meilleurs sur dix}\label{sec:e4}

\paragraph{Question et plan}
Que devient un résultat quand on ne publie que les trois meilleurs de dix essais, une pratique que
\textcite{henderson2018matters} relevaient ? E4 applique cette sélection à chacun des seize couples de E2, puis
compte les comparaisons qui changent de sens parmi les 24 (\cref{sec:plan-e4}).

\paragraph{Une moyenne toujours gonflée}
Le \cref{tab:e4} et la \cref{fig:meilleurs-sur-dix} donnent le gonflement pour chaque couple. Les seize
gonflements sont positifs, par construction. La moyenne gonfle de 32~\% en médiane, de 6~\% à 84~\% selon le
programme et la tâche. Les plus forts touchent PPO JAX sur \task{Walker2d} (84~\%), SAC sur \task{Hopper}
(78~\%) et CrossQ sur \task{HalfCheetah} (74~\%). L'écart type des trois meilleurs vaut en médiane 35~\% de
celui des dix : la sélection gonfle la moyenne et resserre la dispersion apparente, ce qui fait paraître le
résultat plus sûr.

\fig{meilleurs-sur-dix}{Hausse de la moyenne quand on ne publie que les trois meilleurs des dix essais, pour
chaque programme (une couleur) et chaque tâche (une ligne), en pourcentage de la moyenne des dix. La ligne
pointillée marque la médiane, +32~\%. Les seize hausses sont positives : le tri sert toujours celui qui le
pratique. Données : dix graines aléatoires par couple, scores d'évaluation finaux.}

\paragraph{Trois comparaisons inversées}
Sur 24~comparaisons, 3 changent de sens quand chaque programme ne garde que ses trois meilleurs essais :
PPO PyTorch contre PPO JAX sur \task{HalfCheetah} et sur \task{Hopper}, CrossQ contre PPO PyTorch sur
\task{Swimmer}. Deux des trois opposent les deux versions de PPO. Quand l'écart est faible, garder ses
meilleurs essais suffit à changer de vainqueur.

\gentable{e4}

\paragraph{Réplication sur les graines aléatoires 10 à 19}
Refaite à l'identique sur dix autres graines aléatoires, la sélection gonfle la moyenne de 43~\% en médiane, de
7~\% à 85~\%. Elle inverse encore 3~comparaisons sur 24 : CrossQ contre PPO PyTorch sur \task{Hopper}, PPO
PyTorch contre PPO JAX sur \task{Hopper}, CrossQ contre PPO PyTorch sur \task{Swimmer}. Deux des trois
inversions sont les mêmes que sur les graines aléatoires 0 à 9.

\paragraph{Ce que E4 établit}
La sélection des meilleurs essais gonfle toujours la moyenne, d'un tiers environ en médiane. Elle inverse les
comparaisons serrées. Le résultat se réplique sur des graines aléatoires indépendantes. E4 mesure l'effet d'une pratique.
Il ne dit pas combien d'articles la suivent : le sondage mesure cette part (critère X3, \cref{chap:04-sondage-resultats}).

\section{E5 : combien de graines aléatoires ?}\label{sec:e5}

\paragraph{Question et plan}
Combien de graines aléatoires faut-il pour retrouver le bon ordre de deux programmes ? E5 tire avec remise
$n$ essais par programme parmi ceux observés, \num{10000} fois, puis compte la part des tirages qui retrouvent
l'ordre des moyennes de tous les essais (\cref{eq:e5}, \cref{sec:plan-e5}).

\paragraph{Dix graines aléatoires}
Sur les 24 comparaisons, la part médiane des tirages qui retrouvent le bon ordre vaut 94~\% avec trois
graines aléatoires, 97~\% avec cinq et 99~\% avec dix (\cref{tab:e5}). Ce chiffre rassurant tient aux écarts larges entre
nos programmes. Sept comparaisons sont proches, avec des moyennes qui diffèrent de moins de 20~\%. Pour
elles, la médiane tombe à 65~\% avec trois graines aléatoires, 71~\% avec cinq, 77~\% avec dix et 85~\% avec vingt. La
comparaison la plus difficile, CrossQ contre SAC sur \task{HalfCheetah}, ne retrouve le bon ordre que dans
60~\% des tirages avec cinq graines aléatoires.

\gentable{e5}

\paragraph{Vingt graines aléatoires}
Le bon ordre devient celui des moyennes de vingt essais. Les tirages se font parmi les vingt. Sur les 24
comparaisons, la médiane ne bouge presque pas : 93~\% avec trois graines aléatoires, 97~\% avec cinq. Les comparaisons
proches sont désormais huit, plus difficiles que dix graines aléatoires ne le laissaient voir : la médiane vaut 59~\%,
62~\%, 67~\% et 74~\% avec trois, cinq, dix et vingt graines aléatoires. CrossQ et PPO PyTorch sur \task{Swimmer} ont des
moyennes qui diffèrent de moins de 1~\% : même vingt graines aléatoires ne les ordonnent que dans 53~\% des tirages,
autant qu'une pièce de monnaie.

\fig{e5-graines}{Part des \num{10000} tirages qui retrouvent le bon ordre de deux programmes, en fonction du
nombre de graines aléatoires tirées pour chacun (3, 5, 10, 20). Le bon ordre est celui des moyennes de tous les
essais observés. Une courbe par comparaison entre deux programmes sur une tâche, 24 en tout. En rouge, les
comparaisons proches, dont les moyennes diffèrent de moins de 20~\%. Deux panneaux : le bon ordre et les
tirages viennent de dix graines aléatoires, puis de vingt. Avec vingt graines aléatoires comme référence, cinq ne
retrouvent le bon ordre des comparaisons proches que dans 62~\% des tirages en médiane.}

\paragraph{Ce que E5 établit, ses limites}
Cinq graines aléatoires, la valeur médiane que le sondage relève dans les articles de 2026
(\cref{chap:04-sondage-resultats}), suffisent pour des écarts larges. Elles ne suffisent pas pour des écarts de
moins de 20~\% : avec vingt graines aléatoires comme référence, un tel écart se lit à l'envers plus d'une fois sur trois.
E5 fait tenir aux essais observés le rôle de la vraie distribution. Avec dix ou vingt essais, cette
distribution est mal connue : pour les comparaisons les plus proches, le bon ordre l'est aussi.

\section{E6 : une méthode de référence un peu trop belle}\label{sec:e6}

\paragraph{Question et plan}
E6 voulait illustrer la faille F7 : l'écart qu'une nouvelle méthode s'attribue quand elle est seule à être
réglée. Elle oppose, sur \task{Hopper-v5}, dix essais de PPO PyTorch avec les valeurs par défaut de
Stable-Baselines3 à dix essais avec les réglages publiés par RL Baselines3 Zoo, qui font foi pour cette
bibliothèque (\cref{sec:plan-e6}). Le \cref{tab:e6} rassemble tous les groupes.

\paragraph{Premier résultat : les réglages publiés feraient trois fois moins bien}
La version par défaut atteint une IQM de \num{2877} $[\num{2268}, \num{3369}]$. La version réglée atteint
\num{1003} $[856, \num{1146}]$. En moyenne, \num{2776} contre 994, avec un plafond près de \num{1000} atteint
dès \num{50000}~pas. Les réglages publiés feraient près de trois fois moins bien que les valeurs par défaut.

\paragraph{Une erreur de notre côté, qui n'explique rien}
Le résultat était trop beau pour être cru sans vérification. Nous avons trouvé une erreur dans notre code. La
normalisation des récompenses du Zoo emploie le facteur d'actualisation de l'algorithme, 0,999. Notre premier
passage lui laissait la valeur par défaut, 0,99 (écart~2). Le diagnostic, décidé après ce premier résultat,
relance la recette corrigée sur \task{Hopper-v5} et sur \task{Hopper-v4}, plus les valeurs par défaut sur v4,
dix graines aléatoires chacun.
\begin{itemize}
  \item La recette corrigée obtient \num{1012} $[902, \num{1135}]$ sur v5, \num{984} $[773, \num{1002}]$ sur
    v4, la version pour laquelle le recueil a été réglé.
  \item Les valeurs par défaut obtiennent \num{2423} $[\num{1657}, \num{3058}]$ sur v4.
  \item La probabilité qu'un essai corrigé batte un essai avec l'erreur vaut 0,47, de 0,21 à 0,73 :
    l'erreur ne change rien de détectable.
  \item Sur v4, un essai du Zoo corrigé ne bat un essai par défaut qu'avec une probabilité de 0,07, de 0 à
    0,24. L'écart résiste à la correction de Holm ($p=0{,}001$). Sur v5, la même probabilité vaut 0,01.
\end{itemize}
Notre erreur était réelle, sans expliquer l'écart. Elle n'explique pas non plus la différence de version :
sur v4 comme sur v5, la recette plafonne près de \num{1000}.

\paragraph{Le script officiel plafonne aussi}
Deux explications restaient possibles : la recette publiée fait vraiment moins bien que les valeurs par défaut
avec les bibliothèques actuelles, ou notre façon de l'appliquer diffère encore de celle du recueil. L'écart~5,
écrit avant tout entraînement officiel, tranche par le script du recueil lui-même. Les dix politiques
entraînées par \texttt{rl\_zoo3.train} sur \task{Hopper-v4}, évaluées avec notre protocole, obtiennent
\num{995} $[821, \num{1006}]$. C'est autant que notre recette corrigée, \num{984} $[777, \num{1002}]$ dans
cette analyse : la probabilité qu'une politique officielle batte la nôtre vaut 0,58, de 0,32 à 0,82. Les
valeurs par défaut obtiennent \num{2423} $[\num{1687}, \num{3058}]$. Une politique officielle ne bat une
politique par défaut qu'avec une probabilité de 0,09, de 0 à 0,30. L'écart résiste à la correction de Holm
($p=0{,}001$). Les intervalles de \num{984} et de \num{2423} diffèrent de quelques points entre les deux
analyses, parce que chacune tire ses propres rééchantillons.

\gentable{e6}

\paragraph{Ce que E6 établit}
Selon la règle de lecture écrite d'avance, notre application de la recette n'était pas en cause. Avec
Stable-Baselines3~2.9.2, Gymnasium~1.1.1 et MuJoCo~4.0.0, la recette réglée et publiée par le recueil fait
moins bien que les valeurs par défaut sur \task{Hopper}, sur sa version 4 comme sur la version 5. E6 ne montre
donc pas une référence mal réglée par ses auteurs. Il montre qu'un réglage publié fait moins bien, aujourd'hui,
que les valeurs par défaut de la même bibliothèque. L'explication la plus probable est que ce réglage a vieilli
avec les logiciels qui l'entourent. E6 rappelle aussi qu'un résultat trop beau se vérifie jusqu'au script
officiel.

\paragraph{Limites}
Nous ne savons pas quel changement de bibliothèque a fait vieillir la recette. Nous n'avons essayé qu'une
tâche et qu'un algorithme du recueil. Le plafond près de \num{1000} correspond à peu près à ce que rapporte, sur
\num{1000}~pas, la prime donnée à un sauteur qui reste debout : il évoque un agent qui tient debout sans
avancer. Nous ne l'avons pas vérifié sur les vidéos de chaque politique.

\section{E7 : ce que coûte une physique modifiée}\label{sec:e7}

\paragraph{Question et plan}
Pineau demandait de séparer l'entraînement de l'évaluation. Nos agents sont déjà évalués sur des graines
aléatoires jamais vues à l'entraînement. E7 va plus loin : chaque politique finale joue aussi dans une
physique modifiée, avec des masses alourdies de 20~\% et des frottements de glissement réduits de 20~\%. Pour
chaque essai, on calcule la part du score gardée, le rapport du score en physique modifiée au score en physique
d'entraînement, puis son IQM sur les dix graines aléatoires avec son intervalle (\cref{sec:plan-e7}).

\paragraph{Douze couples sur seize perdent}
Le \cref{tab:e7} et la \cref{fig:e7-carte} donnent la part gardée pour chaque couple. Douze couples sur seize
perdent une part établie de leur score, c'est-à-dire avec un intervalle qui reste sous 1. Trois ne changent pas
de façon détectable. Un gagne : SAC sur \task{Swimmer}, à 1,07 $[1{,}02, 1{,}12]$. La perte dépend du
programme et de la tâche.
\begin{itemize}
  \item Sur \task{Hopper}, PPO ne garde que 22~\% de son score dans sa version PyTorch et 18~\% dans sa
    version JAX, quand SAC en garde 64~\% et CrossQ 65~\%.
  \item Sur \task{Walker2d}, PPO PyTorch garde 34~\%, CrossQ 46~\%, PPO JAX 52~\%, SAC 84~\%.
  \item Sur \task{HalfCheetah}, les pertes sont plus modestes : de 85~\% gardés pour PPO PyTorch à 92~\% pour
    PPO JAX.
  \item Sur \task{Swimmer}, aucun programme ne perd.
\end{itemize}

\fig{e7-carte}{Part du score gardée dans une physique modifiée (masses multipliées par 1,2, frottements de
glissement par 0,8), pour chaque programme et chaque tâche. Chaque case donne la
moyenne interquartile, sur dix graines aléatoires, du rapport entre le score dans la physique modifiée et le
score dans la physique d'entraînement, les deux mesurés sur les mêmes graines aléatoires d'évaluation. Les tâches sont en
colonnes, les programmes en lignes. Une case marquée d'une perte a un intervalle à 95~\% qui reste sous 1, une case marquée d'un gain un intervalle au-dessus de
1.}

\gentable{e7}

\paragraph{Le classement bouleversé sur Hopper}
Le complément descriptif de l'écart~11, sans test, classe les programmes dans chaque physique par IQM du
score. Sur \task{Hopper}, l'ordre se renverse. PPO JAX, premier dans la physique d'entraînement avec
\num{3064}, finit dernier dans la physique modifiée avec 525. CrossQ, troisième avec \num{2302}, prend la tête
avec \num{1302}, devant SAC (751) et PPO PyTorch (609). Sur \task{Walker2d}, SAC passe de la dernière place à
la deuxième. Sur \task{HalfCheetah} et \task{Swimmer}, l'ordre ne change pas. La \cref{fig:physique-modifiee}
montre le renversement de \task{Hopper}.

\fig{physique-modifiee}{\task{Hopper}, quatre programmes, dix graines aléatoires chacun. À gauche, le score dans
la physique de l'entraînement. À droite, le score avec des masses alourdies de 20~\% et des frottements
réduits de 20~\%. Chaque ligne relie le score et le rang d'un programme dans les deux physiques. PPO écrit en
JAX passe de la première à la dernière place, CrossQ de la troisième à la première. Les scores sont des
moyennes interquartiles des scores d'évaluation finaux. Les rangs les comparent sans test.}

\paragraph{Vingt graines aléatoires}
Les verdicts restent les mêmes pour 14~couples sur 16. Deux changent. SAC sur \task{Walker2d} ne perd plus de
façon détectable : 0,91, de 0,79 à 1,01. CrossQ sur \task{Swimmer} gagne un peu : 1,03, de 1,01 à 1,04. Il y a
ainsi 11~pertes et 2~gains. Le classement descriptif se renverse toujours sur \task{Hopper}, seule tâche où le
programme en tête change.

\paragraph{Ce que E7 établit, ses limites}
Un programme en tête dans la physique d'entraînement finit dernier dans une physique voisine : PPO JAX sur
\task{Hopper}. Le classement
établi dans le monde de l'entraînement dit peu de chose de la robustesse. E7 ne teste qu'une modification,
qui change deux grandeurs à la fois. La section suivante sépare les deux grandeurs et fait varier l'ampleur du
changement. Les classements dans chaque physique comparent des IQM sans test. La part gardée se calcule essai
par essai : elle ne se lit pas comme le rapport des deux IQM du classement.

\section{La physique modifiée par degrés}\label{sec:dose}

\paragraph{Question et plan}
Le renversement de E7 tient-il à une seule modification, choisie d'avance ? L'écart~12, point~ii, réévalue
chaque politique finale en ne changeant qu'une grandeur à la fois : les masses de 30~\% plus légères à 30~\% plus
lourdes, les frottements de 40~\% plus faibles à 20~\% plus forts. Un couple perd quand l'intervalle de l'IQM
de son rapport au score d'origine reste sous 1, comme dans E7 (\cref{sec:plan-e7}).

\paragraph{Le contrôle}
À 1,0 fois, les réévaluations retrouvent le score final écrit par \file{run.py} : le plus grand écart
relatif vaut zéro, pour les 160 politiques à dix graines aléatoires comme pour les 320 à vingt.

\paragraph{La perte croît avec l'ampleur du changement}
Le \cref{tab:dose}, la \cref{fig:intensites} et la \cref{fig:dose-carte} comptent les couples qui perdent, sur
16, pour chaque physique. À dix graines aléatoires :
\begin{itemize}
  \item masses multipliées par 0,7, 0,8, 0,9 : 8, 6 et 3~couples perdent ;
  \item masses multipliées par 1,1, 1,2, 1,3 : 9, 12 et 11~couples perdent ;
  \item frottements multipliés par 0,6, 0,7, 0,8, 0,9 : 7, 6, 4 et 3~couples perdent ;
  \item frottements multipliés par 1,1 et 1,2 : 1 et 9~couples perdent.
\end{itemize}
Les pertes dessinent un V autour de la physique d'origine, plus raide du côté des masses alourdies. Alourdir les
masses de 20~\% fait perdre à lui seul 12~couples, autant que la modification combinée de E7. Réduire les
frottements de 20~\% n'en fait perdre que 4. Dans E7, la perte semble venir surtout des masses, même si les deux
expériences ne testent pas cette décomposition. Quelques couples gagnent, surtout quand les frottements
baissent : 3 gains à 0,6 fois et à 0,7 fois.

\gentable{dose}

\fig{intensites}{Quatre programmes sur quatre tâches, dix graines aléatoires chacun, soit 16 couples de
programme et de tâche. Chaque point compte les couples qui perdent une part détectable de leur score quand une
seule grandeur de la physique change : les masses (de 30~\% plus légères à 30~\% plus lourdes) ou les
frottements de glissement (de 40~\% plus faibles à 20~\% plus forts). Un couple perd quand l'intervalle à
95~\% de son score, rapporté à celui de la physique d'entraînement, reste sous 1. À 0~\%, la physique
d'entraînement sert de référence : aucune perte n'y est possible.}

\fig{dose-carte}{Part du score gardée dans douze physiques modifiées, avec dix graines aléatoires puis vingt,
côte à côte. Une rangée par programme et par tâche (16), une colonne par physique : masses de 70~\% à 130~\%
de leur valeur d'entraînement, puis frottements de glissement de 60~\% à 120~\%, une seule grandeur changée à
la fois. Une case n'est colorée que si l'intervalle à 95~\% de la moyenne interquartile du rapport au score
d'origine évite 1 : en rouge une perte, plus foncée quand plus de la moitié du score est perdue, en vert un
gain. Sous chaque colonne, le nombre de couples qui perdent, sur 16.}

\paragraph{Le programme en tête de Hopper selon la physique}
Sur \task{Hopper}, PPO JAX, premier dans la physique d'origine, n'est plus premier dans 10 des 12 physiques
modifiées. Il ne garde la tête qu'avec des frottements à 0,9 et 1,1 fois. La tête passe à PPO PyTorch ou à
CrossQ selon la grandeur et le facteur. Sur les trois autres tâches, le programme en tête change peu : SAC dans
les douze physiques de \task{HalfCheetah}, CrossQ dans les douze de \task{Walker2d}, PPO JAX dans dix des douze
de \task{Swimmer}.

\paragraph{Vingt graines aléatoires}
Le V garde sa forme. Avec des masses à 0,9, 1,1, 1,2 et 1,3 fois, 4, 11, 12 et 12~couples perdent. Avec des
frottements à 1,2 fois, 10. Sur \task{Hopper}, PPO JAX perd toujours la tête dans 10 physiques sur 12. Sur
\task{HalfCheetah}, CrossQ prend la tête dans 9 des 12 physiques, comme il la prend dans la physique d'origine
à vingt graines aléatoires.

\paragraph{Ce que la physique par degrés établit}
Le renversement de \task{Hopper} ne tient pas à une modification choisie. Il apparaît dans presque toutes les
physiques voisines, à dix comme à vingt graines aléatoires. Les classements par physique restent sans test. Les
modifications restent petites et uniformes : chaque masse ou chaque frottement change du même facteur. Elles ne
remplacent pas un vrai test de généralisation, sur d'autres corps ou d'autres tâches.

\section{B3 : la mesure change-t-elle le classement ?}\label{sec:b3}

\paragraph{Question et plan}
Un article qui ne définit pas sa mesure peut-il choisir son vainqueur ? B3 classe les quatre programmes selon
trois mesures calculées sur les mêmes courbes : le score final, le meilleur score atteint le long de la courbe
et la moyenne de la courbe (\cref{sec:plan-b3}).

\paragraph{Dix graines aléatoires : huit inversions, aucune établie}
Huit paires de programmes changent d'ordre d'une mesure à l'autre (\cref{tab:b3}). Le programme en tête change
de mesure en mesure sur deux tâches.
\begin{itemize}
  \item Sur \task{HalfCheetah}, la moyenne de la courbe place CrossQ devant SAC, en tête selon le score final
    et le meilleur score.
  \item Sur \task{Hopper}, le meilleur score place PPO PyTorch devant PPO JAX, en tête selon les deux autres
    mesures.
\end{itemize}
Les autres inversions touchent CrossQ et PPO PyTorch sur \task{Hopper} et sur \task{Swimmer}. Aucune n'est
établie : elles opposent des programmes proches, que dix graines aléatoires ne départagent sous aucune mesure. La moyenne
de la courbe couvre un million de pas pour PPO et \num{200000} pour SAC et CrossQ. Elle favorise un programme
qui apprend vite au début.

\gentable{b3}

\paragraph{Vingt graines aléatoires et budget égal}
\begin{itemize}
  \item À budgets inégaux et vingt graines aléatoires, 8~paires changent encore d'ordre, sans inversion établie. Le
    programme en tête ne change de mesure en mesure que sur \task{Hopper}.
  \item À budget égal et dix graines aléatoires, 4~inversions, aucune établie. Le programme en tête ne change de mesure
    en mesure que sur \task{HalfCheetah}, entre SAC et CrossQ.
  \item À budget égal et vingt graines aléatoires, 6~inversions, dont 1 établie : sur \task{Swimmer}, PPO PyTorch et SAC
    changent d'ordre entre le meilleur score atteint et la moyenne de la courbe. Le programme en tête ne change
    de mesure en mesure sur aucune tâche.
\end{itemize}

\paragraph{Ce que B3 établit}
Entre deux programmes proches, la mesure choisie décide parfois du vainqueur apparent. Entre programmes
nettement séparés, elle ne change rien. Une seule inversion est établie dans les quatre analyses. Elle
oppose deux programmes loin du premier. Un article qui ne dit pas sa mesure laisse ouverte, pour les
comparaisons serrées, la question du vainqueur.

\section{Dix ou vingt graines aléatoires, budget égal : ce qui tient et ce qui change}\label{sec:consolidation}

Les résultats à dix graines aléatoires restent les résultats principaux, comme le protocole l'annonçait. Les
ajouts de l'écart~12 les confirment ou les nuancent.

\paragraph{Ce qui tient}
\begin{itemize}
  \item Le programme en tête change avec la tâche aux budgets du protocole (E2), de façon établie, avec ou sans
    correction de Bonferroni.
  \item La sélection des trois meilleurs essais sur dix gonfle toujours la moyenne et inverse 3~comparaisons
    sur 24, sur les graines aléatoires 10 à 19 comme sur les graines aléatoires 0 à 9 (E4).
  \item La physique modifiée renverse le classement de \task{Hopper}. Les verdicts de E7 restent les mêmes
    pour 14~couples sur 16. La physique par degrés garde sa forme en V.
  \item Les inversions de B3 à budgets inégaux restent au nombre de 8, sans inversion établie.
\end{itemize}

\paragraph{Ce qui change}
\begin{itemize}
  \item Sur \task{HalfCheetah}, CrossQ prend la tête à SAC, sans écart établi entre eux (E2).
  \item Les deux versions de PPO se séparent sur \task{Walker2d} et \task{Swimmer}, dans des sens opposés (E3).
  \item Les comparaisons proches sont plus difficiles à ordonner : avec cinq graines aléatoires, le bon ordre n'est
    retrouvé que dans 62~\% des tirages en médiane, au lieu de 71~\% (E5).
  \item À budget égal et vingt graines aléatoires, CrossQ mène sur les quatre tâches : le classement ne change plus avec la
    tâche (E2).
\end{itemize}
Le dernier point est la nuance qui pèse le plus. Le verdict de E2, tel qu'il était annoncé, tient. Mais à
vingt graines aléatoires, le changement de programme en tête d'une tâche à l'autre s'explique en bonne part par le budget
cinq fois plus grand de PPO.

\section{Ce que ces expériences ne disent pas}\label{sec:limites-experiences}

Les limites propres à chaque expérience sont dites dans sa section. Quelques limites valent pour toutes.
\begin{itemize}
  \item Tous les essais ont tourné sur une seule machine, un Mac à puce Apple M2 Max, sur processeur. Une
    relance ailleurs donnerait d'autres courbes pour les mêmes graines aléatoires.
  \item Les programmes gardent les valeurs par défaut de leurs bibliothèques. Les classements valent pour ces
    réglages. Ils ne disent rien des algorithmes au mieux de leur forme.
  \item Quatre tâches MuJoCo, toutes des corps plans ou presque, ne représentent pas l'apprentissage par
    renforcement de 2026. Le volet des modèles de langue (\cref{chap:07-modeles-de-langue}) en touche une autre
    part, à petite échelle.
  \item Dix épisodes d'évaluation, joués sans hasard sur dix positions de départ fixes, mesurent une politique
    sur une petite partie de ses situations possibles.
  \item Dix graines aléatoires, puis vingt, laissent des intervalles larges pour les comparaisons serrées. E5
    le chiffre.
  \item Plusieurs analyses ont été ajoutées après avoir vu un résultat : la seconde statistique de E1, le
    diagnostic de E6, le classement descriptif de E7. La \cref{sec:avant-apres} les liste. Le texte les
    présente comme telles.
\end{itemize}
